% Front matter: portada, segunda portada, créditos, autores, agradecimientos
% Mismo estilo que Financieras: página 1 = imagen de fondo a sangre (luego se insertan logo y título);
% página 2 = fondo blanco, título, autores, logo Ediciones CVEA, EECA/UCV.

\pagestyle{empty}

% ----- Página 1: Portada (imagen compuesta: fondo verde + logo horizontal grande + Ediciones CVEA) -----
\newgeometry{left=0pt,right=0pt,top=0pt,bottom=0pt}
\setlength{\parindent}{0pt}
\setlength{\parskip}{0pt}
\noindent
\IfFileExists{portada-final.png}{%
  \includegraphics[width=1.02\paperwidth,height=1.02\paperheight,keepaspectratio]{portada-final.png}%
}{\IfFileExists{portada-actuariales-propuesta.png}{%
  \includegraphics[width=1.02\paperwidth,height=1.02\paperheight,keepaspectratio]{portada-actuariales-propuesta.png}%
}{%
  \begin{center}
  \vspace*{3cm}
  {\Huge\bfseries Introducción a las Matemáticas Actuariales}\\[1.5em]
  {\Large Material docente --- Matemáticas Actuariales I}\\[2cm]
  {\large Prof. Felipe Moreno \quad Prof. Jorge Días \quad Prof. Angel Colmenares}\\[1em]
  {\large CVEA / EECA-UCV}
  \end{center}
}}
\newpage
\restoregeometry

% ----- Página 2: Segunda portada (interior, fondo blanco) -----
\thispagestyle{empty}
\vspace*{1.2cm}
\begin{center}
{\Huge\bfseries Introducción}\\[0.6em]
{\Huge\bfseries a las Matemáticas Actuariales}\\[1.5em]
{\Large Material docente --- Matemáticas Actuariales I}\\[1.8cm]
{\Large
Prof. Felipe Moreno\\
Prof. Jorge Días\\
Prof. Angel Colmenares}\\[2.2cm]
\IfFileExists{assets/logos/Ediciones_CVEA_Logo_BN.png}{%
  \includegraphics[width=0.5\textwidth,keepaspectratio]{assets/logos/Ediciones_CVEA_Logo_BN.png}\\[1.2cm]
}{\IfFileExists{Ediciones_CVEA_Logo_BN.png}{%
  \includegraphics[width=0.5\textwidth,keepaspectratio]{Ediciones_CVEA_Logo_BN.png}\\[1.2cm]
}{}}
{\Huge\bfseries Ediciones CVEA}
\end{center}
\vfill
\begin{center}
{\large
Escuela de Estadística y Ciencias Actuariales (EECA)\\[0.4em]
Universidad Central de Venezuela (UCV)}
\end{center}
\vspace{1.5cm}
\newpage

% ----- Página 3: Créditos / Derechos reservados -----
\thispagestyle{empty}
\vspace*{1.5cm}
\begin{center}
\textbf{Introducción a las Matemáticas Actuariales}\\
Material docente --- Matemáticas Actuariales I\\[1cm]
Primera edición\\[2cm]
\end{center}

Prohibida la reproducción total o parcial de esta obra, por cualquier medio, sin la autorización escrita del editor.

\vspace{1em}
\textbf{DERECHOS RESERVADOS} \copyright\ \the\year, respecto de la primera edición por\\
\textbf{Ediciones CVEA}\\
Centro Venezolano de Estudios Actuariales (CVEA)\\
Escuela de Estadística y Ciencias Actuariales (EECA)\\
Universidad Central de Venezuela

\vspace{0.8em}
Av. Los Ilustres, Los Chaguaramos\\
Caracas 1040, Venezuela

\vspace{0.5em}
\textit{Primera edición}\\
ISBN: 978-XXX-XXX-XXX-X (por asignar)

\vspace{1em}
\small
Impreso en Venezuela \textit{/} Printed in Venezuela

\vfill
\newpage

% ----- Página 4: Sobre esta obra (tributo y reconocimientos) -----
\thispagestyle{empty}
\vspace*{1.2cm}
\begin{center}
{\LARGE\bfseries Sobre esta obra}
\end{center}
\vspace{1.2cm}

Este material es la \textbf{Introducción a las Matemáticas Actuariales} (Matemáticas Actuariales I), desarrollada en el marco del Centro Venezolano de Estudios Actuariales (CVEA) y la Escuela de Estadística y Ciencias Actuariales (EECA) de la Universidad Central de Venezuela. El contenido base proviene del curso \textit{Matemáticas Actuariales I} (período lectivo II-2024). Se asume como prerrequisito el libro \textit{Introducción a las Matemáticas Financieras} (Ediciones CVEA).

\vspace{1em}
\textbf{Reconocimiento al Prof. Felipe Moreno.} Este libro es un reconocimiento y un tributo al \textbf{Profesor Felipe Moreno}, por parte del \textbf{Centro Venezolano de Estudios Actuariales (CVEA)} y de los profesores \textbf{Jorge Días} y \textbf{Angel Colmenares}, compiladores de este material y en su momento alumnos del Prof. Moreno. El Prof. Felipe Moreno se graduó como actuario en 1990 y poco después inició su labor docente en la EECA-UCV. A lo largo de más de tres décadas ha formado a generaciones de actuarios con una exigencia y un rigor en la materia que siempre lo han caracterizado. Hoy, próximo a su jubilación, este texto busca recoger y preservar parte de su conocimiento y de su manera de enseñar las matemáticas actuariales: la claridad en los conceptos, el encadenamiento lógico de los temas y la alta exigencia que ha marcado a sus alumnos. Con esta obra, el CVEA y los profesores Días y Colmenares le expresan su gratitud y su reconocimiento por tantos años de intensa labor al servicio de la formación actuarial en Venezuela.

\vfill
\newpage

% ----- Página 5: Agradecimientos -----
\thispagestyle{empty}
\vspace*{1.5cm}
\begin{center}
{\large\bfseries Agradecimientos}
\end{center}
\vspace{1.5cm}

El CVEA agradece a la Escuela de Estadística y Ciencias Actuariales (EECA) de la Universidad Central de Venezuela por el apoyo institucional para la elaboración de este material docente y a los profesores y estudiantes cuyas observaciones enriquecen el proyecto.

\vfill
\newpage
\pagestyle{plain}
\let\cleardoublepage\savedcleardoublepage
